% Figura 2.3 — Esquema conceptual de una red generativa antagónica condicional de Wasserstein
% con crítico de parches y término de reconstrucción (esquema sin datos).
% Fuente editable (TikZ). Compilar con:  pdflatex fig2-3_gan_condicional.tex
% Notación igual a la del capítulo 3: entrada u, referencia y, generador G_theta, crítico D_phi
% con puntuaciones por parche y puntuación media S_phi, interpolación y~. Paleta Paul Tol "bright";
% generador y crítico se distinguen también por el trazo del borde (continuo y discontinuo).
\documentclass[12pt,border=2pt]{standalone}
\usepackage[T1]{fontenc}
\usepackage[utf8]{inputenc}
\usepackage[spanish]{babel}
\usepackage{lmodern}
\usepackage{amsmath,amssymb}
\usepackage{bm}
\usepackage{tikz}
\usetikzlibrary{babel,arrows.meta,positioning,calc}

\definecolor{tolblue}{HTML}{4477AA}
\definecolor{tolgreen}{HTML}{228833}
\definecolor{tolpurple}{HTML}{AA3377}
\definecolor{tolgrey}{HTML}{BBBBBB}

\begin{document}
\scriptsize
\hyphenpenalty=10000 \exhyphenpenalty=10000
\begin{tikzpicture}[
  >={Stealth[length=1.8mm]},
  img/.style={draw=black!70, line width=0.5pt, minimum width=10mm, minimum height=10mm, inner sep=0pt, fill=tolgrey!25},
  caja/.style={draw=#1!85!black, fill=#1!10, rounded corners=1.5pt, line width=0.6pt, align=center, inner sep=2pt},
  flecha/.style={->, line width=0.6pt, draw=black!75},
  nota/.style={font=\tiny, text=black!75, align=center}
]
% ---------------- Entrada de baja estadística -----------------------------------
\node[img] (u) at (0,0) {};
\foreach \i in {1,...,28} {\fill[black!70] ($(u.south west)+({mod(\i*37,97)/97*0.9+0.05},{mod(\i*61,89)/89*0.9+0.05})$) circle (0.25pt);}
\node[nota, below=1mm of u] {entrada $\bm u$\\ baja estadística};

% ---------------- Generador U-Net (icono) ---------------------------------------
\begin{scope}[shift={(1.0,-0.55)}]
  \foreach \k/\h in {0/1.1,1/0.8,2/0.5} {\fill[tolblue!25, draw=tolblue!85!black, line width=0.4pt] ({0.3*\k},{0.55-\h/2}) rectangle ({0.3*\k+0.18},{0.55+\h/2});}
  \fill[tolblue!45, draw=tolblue!85!black, line width=0.4pt] (0.9,0.4) rectangle (1.08,0.7);
  \foreach \k/\h in {0/0.5,1/0.8,2/1.1} {\fill[tolblue!25, draw=tolblue!85!black, line width=0.4pt] ({1.2+0.3*\k},{0.55-\h/2}) rectangle ({1.2+0.3*\k+0.18},{0.55+\h/2});}
  \draw[->, tolblue!85!black, line width=0.35pt] (0.09,1.12) -- (0.09,1.24) -- (1.89,1.24) -- (1.89,1.12);
  \draw[->, tolblue!85!black, line width=0.35pt] (0.39,0.97) -- (0.39,1.06) -- (1.59,1.06) -- (1.59,0.97);
  \coordinate (Gin) at (-0.05,0.55); \coordinate (Gout) at (2.12,0.55);
  \node[nota, text=tolblue!80!black, align=center, anchor=south] at (0.99,1.3) {generador $G_\theta$ (U-Net)};
\end{scope}
\draw[flecha] (u.east) -- (Gin);

% ---------------- Estimación y referencia ---------------------------------------
\node[img] (yhat) at (4.1,0) {};
\foreach \r/\c in {0.36/20,0.26/35,0.16/50,0.07/65} {\fill[black!\c] (yhat.center) circle (\r);}
\node[nota, below=1mm of yhat] {estimación\\ $G_\theta(\bm u)$};
\draw[flecha] (Gout) -- (yhat.west);

\node[img] (y) at (4.1,1.75) {};
\foreach \r/\c in {0.36/20,0.26/35,0.16/50,0.07/65} {\fill[black!\c] (y.center) circle (\r);}
\foreach \i in {1,...,10} {\fill[black!70] ($(y.south west)+({mod(\i*41,97)/97*0.9+0.05},{mod(\i*53,89)/89*0.9+0.05})$) circle (0.2pt);}
\node[nota, left=1mm of y] {referencia $\bm y$\\ alta estadística};

% ---------------- Crítico condicional de parches ---------------------------------
\node[caja=tolpurple, dashed, text width=17mm, minimum height=12mm] (crit) at (6.45,0.875)
  {crítico $D_\phi$\\ evalúa el par\\ $(\bm u,\cdot)$};
\draw[flecha] (yhat.east) -- ++(0.5,0) |- ($(crit.west)+(0,-0.3)$);
\draw[flecha] (y.east) -- ++(0.5,0) |- ($(crit.west)+(0,0.3)$);
\draw[flecha, black!55, dash pattern=on 2pt off 1.2pt] (u.north) -- (0,2.75) -| (crit.north);
\node[nota, anchor=south] at (2.2,2.75) {la entrada condiciona al crítico};

\begin{scope}[shift={(7.75,0.45)}]
  \foreach \i in {0,...,4} \foreach \j in {0,...,4} {
    \pgfmathsetmacro{\v}{20+mod(\i*13+\j*29,60)}
    \fill[tolpurple!\v] ({0.17*\i},{0.17*\j}) rectangle ({0.17*\i+0.17},{0.17*\j+0.17});}
  \draw[black!70, line width=0.4pt] (0,0) rectangle (0.85,0.85);
  \draw[black, line width=0.7pt] (0.34,0.34) rectangle (0.51,0.51);
  \node[nota, anchor=north] at (0.425,-0.05) {puntuación\\ por parche};
\end{scope}
\draw[flecha] (crit.east) -- (7.75,0.875);
\node[caja=tolpurple, dashed, text width=12mm] (S) at (9.75,0.875) {media\\ $S_\phi(\bm u,\cdot)$};
\draw[flecha] (8.6,0.875) -- (S.west);

% ---------------- Términos de pérdida y actualización ----------------------------
\node[caja=tolgreen, text width=30mm] (l1) at (4.1,-2.05) {reconstrucción $\mathcal L_1$: discrepancia píxel a píxel entre $G_\theta(\bm u)$ y $\bm y$};
\draw[flecha, tolgreen!70!black] ($(yhat.south)+(0,-0.62)$) -- (l1.north);

\node[caja=tolpurple, dashed, text width=36mm] (adv) at (8.6,-2.05) {estimación de la distancia de Wasserstein y penalización del gradiente en puntos interpolados $\widetilde{\bm y}$};
\draw[flecha, tolpurple!80!black] (S.south) -- (S.south |- adv.north);

\node[caja=tolblue, text width=66mm] (upd) at (6.35,-3.55) {actualización alternada: el crítico aprende a separar pares reales y generados;\\ el generador reduce $\mathcal L_1$ y esa separación};
\draw[flecha] (l1.south) -- (l1.south |- upd.north);
\draw[flecha] (adv.south) -- (adv.south |- upd.north);
\end{tikzpicture}
\end{document}
